%%%%%%%%%%%%%%%%%%%%%%%%%%%%%%%%%%%%%%%%%%%%%%%%%%%%%%%%%%%%%%%%%%%%%%%%

%%% LaTeX Template for ECAI Papers 
%%% Prepared by Ulle Endriss (version 1.0 of 2023-12-10)

%%% To be used with the ECAI class file ecai.cls.
%%% You also will need a bibliography file (such as mybibfile.bib).

%%%%%%%%%%%%%%%%%%%%%%%%%%%%%%%%%%%%%%%%%%%%%%%%%%%%%%%%%%%%%%%%%%%%%%%%

%%% Start your document with the \documentclass{} command.
%%% Use the first variant for the camera-ready paper.
%%% Use the second variant for submission (for double-blind reviewing).

% \documentclass{ecai} 
\documentclass[doubleblind]{ecai} 

%%%%%%%%%%%%%%%%%%%%%%%%%%%%%%%%%%%%%%%%%%%%%%%%%%%%%%%%%%%%%%%%%%%%%%%%

%%% Load any packages you require here. 

\usepackage{latexsym}
\usepackage{amssymb}
\usepackage{amsmath}
\usepackage{amsthm}
\usepackage{booktabs}
\usepackage{enumitem}
\usepackage{graphicx}
\usepackage{subcaption}
\usepackage{multicol}
\usepackage{tabularx}
\usepackage{makecell}
\usepackage[utf8]{inputenc}
\usepackage{multirow}
\usepackage{color}
\usepackage{textcomp}
\usepackage{array}
\usepackage{float} 
\usepackage{arydshln}
% \usepackage[hyphens]{url}
\usepackage{xurl}


%%%%%%%%%%%%%%%%%%%%%%%%%%%%%%%%%%%%%%%%%%%%%%%%%%%%%%%%%%%%%%%%%%%%%%%%

%%% Define any theorem-like environments you require here.

\newtheorem{theorem}{Theorem}
\newtheorem{lemma}[theorem]{Lemma}
\newtheorem{corollary}[theorem]{Corollary}
\newtheorem{proposition}[theorem]{Proposition}
\newtheorem{fact}[theorem]{Fact}
\newtheorem{definition}{Definition}

%%%%%%%%%%%%%%%%%%%%%%%%%%%%%%%%%%%%%%%%%%%%%%%%%%%%%%%%%%%%%%%%%%%%%%%%

%%% Define any new commands you require here.

\newcommand{\BibTeX}{B\kern-.05em{\sc i\kern-.025em b}\kern-.08em\TeX}

%%%%%%%%%%%%%%%%%%%%%%%%%%%%%%%%%%%%%%%%%%%%%%%%%%%%%%%%%%%%%%%%%%%%%%%%

\begin{document}

%%%%%%%%%%%%%%%%%%%%%%%%%%%%%%%%%%%%%%%%%%%%%%%%%%%%%%%%%%%%%%%%%%%%%%%%

\begin{frontmatter}

%%% Use this command to specify your submission number.
%%% In doubleblind mode, it will be printed on the first page.

\paperid{5628} 

%%% Use this command to specify the title of your paper.

% Refine title:
\title{Exploring Structural Brain Connectivity in Term and Preterm Infants with Explainable AI and Fuzzy Logic}
% Needs to mention fuzzy logic in the title

% Ideas for titles:
% 1) Exploring the structural connectivity of term and pre-term infant brains via machine learning explainability

% 2)  

%%% Use this combinations of commands to specify all authors of your 
%%% paper. Use \fnms{} and \snm{} to indicate everyone's first names 
%%% and surname. This will help the publisher with indexing the 
%%% proceedings. Please use a reasonable approximation in case your 
%%% name does not neatly split into "first names" and "surname".
%%% Specifying your ORCID digital identifier is optional. 
%%% Use the \thanks{} command to indicate one or more corresponding 
%%% authors and their email address(es). If so desired, you can specify
%%% author contributions using the \footnote{} command.

\author[A]{\fnms{First}~\snm{Author}\orcid{....-....-....-....}\thanks{Corresponding Author. Email: somename@university.edu.}\footnote{Equal contribution.}}
\author[B]{\fnms{Second}~\snm{Author}\orcid{....-....-....-....}\footnotemark}
\author[B,C]{\fnms{Third}~\snm{Author}\orcid{....-....-....-....}} 

\address[A]{Short Affiliation of First Author}
\address[B]{Short Affiliation of Second Author and Third Author}
\address[C]{Short Alternate Affiliation of Third Author}


\begin{abstract}

%Final Abstract
Preterm births have been associated with altered neurological development for neonatal infants, this has been implicated in certain neurodevelopmental conditions in later life. Advances in brain imaging methods, such as Magnetic Resonance Imaging, have allowed for the analysis of physical connectivity of brain matter in infants shortly after birth. However, commonly used methods of investigating such data rely on a brain network analysis, traditionally based on graph-theoretical approaches, which may fail to capture complex patterns involving both local and global network structures and spatial information. Furthermore, many previous studies of infant brain data rely on a priori selection of specific graph connectivity measures. We propose employing machine learning models such as logistic regression and Graph Neural Networks to provide a data-driven approach for classifying preterm and term brain networks at birth. We utilize fuzzy logic, and explainability methods including Shapley Additive Explanations to identify influential regions and connections in decision making. In our analysis, brain regions are represented as spatially embedded nodes, with edges representing strength of structural connections between areas. Using this setup, our model achieves a binary classification accuracy of 88.57\%. This performance is further enhanced using a fuzzy boundary between preterm and term classes, achieving an accuracy of 96.19\%. This demonstrates that the model can be assisted particularly by adding context to 'near-term' born infant cases. These analyses highlight important connections, and key nodes including deep brain structures which are broadly consistent with biological literature.
\end{abstract}

\end{frontmatter}   



%%%%%%%%%%%%%%%%%%%%%%%%%%%%%%%%%%%%%%%%%%%%%%%%%%%%%%%%%%%%%%%%%%%%%%%%

\section{Introduction}

% CP

% Understanding the potential neurological implications of preterm birth is important. Advances in brain imaging methods has allowed us to visualise the connectivity patterns in infants shortly after their birth. These are usually analysed through graph theoretical measures, and compared statistically in order to observe whether preterm and full-term infants brains have differences in their structure.

Preterm births disrupt the critical final trimester of infants and carry the potential to affect infant neuro-development, potentially leading to further complications later in life \cite{schuetze2016morphological,ball2012effect,ream2018neurologic,lee2021cingulum}. Recent advances in infant brain imaging methods applied shortly after their birth \cite{taoudi2022predicting}, have allowed researchers to map brain matter connectivity into structured network representations suitable for application of computational methods. Traditionally, brain connectivity of infants has been studied using graph theoretical measures, such as the rich-club coefficient and clustering coefficient to expose the structural differences between term and preterm infant brains \cite{taoudi2022predicting,batalle2017early,karolis2016reinforcement}. Various studies using traditional graph theoretical approaches have suggested that preterm infants exhibit reduced connectivity between hub regions in the brain when compared with term infants, specifically during early stages of birth when these measurements are taken \cite{batalle2017early,karolis2016reinforcement,taoudi2022predicting,bauer2017nonlinear}.

%add this main cite: \cite{taoudi2022predicting}

% but this can be problematic as it is not particularly flexible, and you have to choose the measures to use in advance, meaning that it is potentially easy to miss important patterns.

Many of these previous approaches, while providing valuable insights into certain preterm brain connectivity patterns, have inherent limitations. Namely, they require a priori selection of specific graph connectivity measures \cite{cui2022braingb}. The pre-selection of particular hypothesis-driven measures \cite{van2010comparing,cui2022braingb} raises the potential of overlooking complex or unexpected connectivity structures that may be captured using more data-driven computational methods. Furthermore, these graph theoretic measures are sensitive to the methodological choices, such as thresholding, network normalization, and hub definition \cite{van2010comparing}. Combined, these methodological variations contribute to a lack of reproducibility across studies and datasets \cite{van2010comparing}.



% We try to use ML models to investigate how the brain is structured at birth. With the hope to understand if there is a difference fir preterm infants and what that difference is.
% We also try graph NN models to undestand how the neural graph differs and harness their ability to attend to edges and nodes.

%We propose using machine learning for analyzing neonatal brain structures to identify and characterize the differences between preterm and term infants. We further leverage graph neural network (GNN) models <\textcolor{red}{GNN refs here}> to take advantage of their capability of processing both node and edge data within brain networks. Additionally, we explore fuzzy-logic systems <\textcolor{red}{fuzzy logic ref here}> to accommodate the variability around the traditional literature-defined threshold that defines preterm and term infants.

% Most importantly, we want to ensure that we understand what the differences are, between the preterm and full-term infants. Therefore we use a series of explainability methods to visualise and interpret the results.


Beyond the classification of preterm and term infants, in this paper we are primarily interested in identifying the differences in neurological structures that distinguish the two populations. To achieve this, we apply machine learning (ML) explainability techniques that aid in visualizing and understanding the models' decision boundaries \cite{gilpin2018explaining,lundberg2017unified}. Rather than relying solely on machine learning to extract abstract patterns, we adopt a data‑centric approach that iteratively maps model outputs back onto the dataset, and validates neurodevelopmental variance between preterm and term infants, as described in the literature \cite{gilmore2018imaging,ream2018neurologic,batalle2017early,ball2012effect,engle2006recommendation,hou2024sensitivity,jukic2013length,karnati2020late,lee2021cingulum,moutquin2003classification}.
Through our approach that focuses on explainability of the classification task, we hope to enhance the understanding of the neuro-biological differences of term and preterm infants. The findings from this paper have the potential to support future clinical interventions and developmental support strategies.

\begin{itemize}
    \item We present an empirical comparison of different ML and GNN architectures for classifying preterm versus term infants from structural brain data. This provides insights into the abilities of various AI models applied to complex, real-world biological data. 
    \item We investigate the impact of incorporating a biologically informed feature engineering strategy. We use atlas‑based spatial coordinates (centroids) as additional model inputs.
    \item We introduce a novel adaptation of fuzzy logic for label smoothing, specifically designed to address continuous development in medical classification. This approach combats the noisy class boundary and considers domain-specific priors, resulting in improved model accuracy.
    \item We provide insights into model decision-making through the use of SHAP explainability. We do this by aggregating node‑level attributions and edge‑importance matrices, and projecting them onto the brain atlas to visually interpret the models’ decisions. We critically discuss the consistency across different architechtures, and discuss consistency with known neuroscientific literature.
\end{itemize}

% \begin{comment}
% \begin{itemize}
%     \item Applying ML/GAT models to structural brain connectivity data from infants. 
%     \item Using fuzzy logic to aid in classification
%     \item Providing explainability insights into the model
% \end{itemize}
% \end{comment}



\begin{figure*}[h]
    \centering
    \begin{subfigure}[t]{0.45\textwidth}
        \includegraphics[width=\textwidth]{ecai-template/figures/brainconnectivity_new (3) (1).png}
        \label{subfig:brainconnectivity_top}
    \end{subfigure}\hfill
    \begin{subfigure}[t]{0.45\textwidth}
        \includegraphics[width=\textwidth]{ecai-template/figures/brainconnectivity_vague1.png}
        \label{subfig:brainconnectivity_front}
    \end{subfigure}
    \caption{Brain connectivity visualizations from different perspectives. Nodes represent the centroid coordinate of each brain region, and edges denote the connections between them. Here, edges are shown based on their strength (darkness and thickness). Node color scale is showing regions from front-brain (blue) to top-back (green) and lower-back (yellow). These provide coordinate location of nodes relative to each other. \vspace{0.5cm}}
    \label{fig:brainconnectivity_combined}
\end{figure*}

\section{Related work}

Due to the way structural brain data is processed, it often takes the form of a connectivity matrix, with nodes representing brain regions and edges representing connections between them. Graph theoretical measures are commonly used to analyze such brain connectivity data. For example, studies investigate connectivity patterns and population differences by comparing values of graph metrics derived from the structural connectome. In the context of preterm birth, researchers have used measures such as the rich club coefficient, betweenness centrality, and small worldness. These metrics aim to identify important brain regions, developmental patterns, and network efficiency. However, they are often difficult to interpret. When combined, it becomes unclear how much each measure contributes to the observed similarity or difference between networks. For instance, in studies comparing preterm and term infants, findings on the prominence of the rich club are inconsistent. Some report greater \cite{karolis2016reinforcement}, while others report lesser \cite{de2021preterm} rich club organization in preterm infants. These discrepancies may result from differences in normalization, metric definitions, or how multiple measures are integrated. In addition, many of these metrics do not account for the spatial location of brain regions and often require extra graph-based or post hoc analyses to support hypothesis-driven interpretations.

Outside the context of preterm birth, some studies have applied ML models to analyze infant brain data \cite{scheinost2023machine, suranyi2021segmentation}. Support Vector Machine (SVM) was applied to functional brain connectivity data from infants in order to predict between preterm and term \cite{smyser2016prediction}. This study had a low sample size, with only 50 infants, and considered functional rather than structural connectivity. One recent study \cite{taoudi2022predicting} did consider both preterm and term brains; however, it focused on predicting gestational age rather than directly classifying preterm versus term brains. Additionally, the study suffered from a significant lack of data, with only seven infants in each group. The regression model also showed poor performance, with a very low $R^2$ score, and the authors did not report standard ML metrics, making comparisons difficult.

Recently, researchers have proposed that Graph Neural Network (GNN)-based models may be well suited for tasks involving structural brain data. This approach has been applied in some contexts with varying degrees of success. However, differences between preterm and term brains have not been explored using GNNs. Cui et al. \cite{cui2022braingb} proposed using GNNs and highlighted several ways brain connectivity data can be transformed into graph format. However, the reported results showed low accuracy. Messaritaki et al. \cite{messaritaki2019optimization} also explored this idea, demonstrating various methods for defining structural brain data as a graph. They noted that, despite the apparent suitability for GNNs, brain network nodes often lack meaningful values. They emphasized the importance of considering the strength of connections. Importantly, their work did not address differences between preterm and term infant brains.

Medical literature shows that identifying and predicting differences in infants born before their estimated due date can be challenging. This is due to several factors, including the difficulty of accurately estimating due dates and the common assumption that all pregnancies should ideally last the same amount of time. In reality, gestational age (GA) at birth varies widely, with many births occurring within a window around the estimated due date \cite{jukic2013length}. Researchers have noted that assessing the potential risks for preterm infants born near the full-term threshold is particularly difficult \cite{moutquin2003classification, engle2006recommendation, karnati2020late}. In classifying stages of prematurity, fuzzy logic has been proposed in both the biological domain \cite{vlamou2019fuzzy, pappas2006biological, cao2024fuzzy} and neuroscience \cite{godil2011fuzzy, munir2019neuroscience}. However, this approach has not yet been applied in the context of infant brain network connectivity.

% ML and AI have been used for networks derived from brain scans. However, its use in classifying between preterm and full-term brains has been scarce. BrainGB was proposed by \cite{cui2022braingb} using GNNs for a variety of brain network analysis contexts. They highlight a number of different ways that the neural data can be transformed into a format to be used in a GNN. However, results from this method are low. Importantly they do highlight the importance of considering edge weights, however this method has not been applied to premature vs term brains as far as we are aware. Moreover, it has been suggested that despite appearing suitable for GNNs, the nodes in the brain network generally lack meaningful values and are instead numbered numerically. This idea was explored by \cite{messaritaki2019optimization} who also use structural brain data. They do not consider the problem of neural differences between preterm and term brains, nor do they consider GNNs. However, they discuss the impact of different methods of defining the graphs. 
% \cite{TAOUDIBENCHEKROUN2022119319} uses a random forest approach and a deep neural network in differentiating between neural connectomes based on number of weeks gestation. Due to the low sample, each GA category has a very small number of subjects (7). Both methods are moderately successful, particularly the DNN, which outperforms the RF on most metrics. However, there is a large amount of variance that these methods are unable to account for, they also do not report performance metrics such as f1, accuracy, etc, likely due to this being a regression task rather than classification into age groups.

% Within the context of understanding differences between neural connectomes for preterm and term infants, graph theoretical measures are generally used. These can include Rich Club, betweenness-centrality, small worldedness and more. This can then give an indication into the way the networks self-organise. However, these can be hard to interpret. \cite{} outline a series of complex network metrics and compare results from them, however this is a process whereby it is difficult to know the relative importance of each metric in calculating similarity scores between networks. \cite{TAOUDIBENCHEKROUN2022119319} consider the importance of specific connections and regions, using sensitivity analysis on their RF and DNN models, and highlight that there are certain regions across the models that are more impactful when identifying the gestational age of the infants.




% Moreover, recent suggestions of using fuzzy logic in classification tasks in the biological domain indicate the importance with regards to explainability \cite{cao2024fuzzy} \cite{pappas2006biological}. Difficulties in understanding potential risk for premature infants born close to the 'full-term' cut off, have been identified \cite{moutquin2003classification}\cite{engle2006recommendation}\cite{karnati2020late}. However, to the best of our knowledge this is not considered within the context of brain network connectivity. 



% \begin{itemize}
%     \item neuroscience papers using ML/AI for the structural connectivity particularly the predicting risk paper
%     - eg BrainGB - uses graph model 
%         - they note that the brain graphs are different due to lack of inherent meaning in the nodes (ie they are just labeled numerically) 
%             - but we add coordinate information
%         - they outline a number of different ways to define the graphs, and discuss a bit about the implications
%             - they do run a variety of different versions 
%         - they also use MLP for downstream tasks
%         - they do not appear to have particularly good results
%         - they do not use preterm term classification
%     - Eirini Messaritaki paper
%         - considering again how the graphs can be constructed for optimal use 
%         - they consider graph theory
%         - they use structural brain data
%             - but not preterm
%         - they discuss how different processing methods of the raw data can influence how reproducable the graph theoretical metric results are (i.e. difficult to compare between datasets and processing pipelines)
        
        
%     \item neuro papers not using ML/AI for structural connectivity differences in preterm and term
%     \item papers with similar methods for other biological concepts
%     \item papers using fuzzy logic, as well as graph methods for the classification
%     - eg Jin Cao 
%         - fuzzy logic for disease detection
%         - they discuss different explainability and interoperability methods in AI - particularly in health data, and they particularly focus on the fuzzy logic, as being more human understandable and intuitive.
%             - they particularly focus on interpretability with fuzzy logic
% \end{itemize}
%%%%%%%%%%%%%%%%%%%%%%%%%%%%%%%%%%%%%%%%%%%%%%%%%%%%%%%%%%%%%%%%%%%%%%%%

\section{Method}

We frame the task of distinguishing preterm from term infant brains as a binary classification problem using structural connectomes derived from Diffusion Tensor Imaging (DTI). Each subject’s brain is represented by a connectivity matrix $A_i \in \mathbb{R}^{n \times n}$, where the rows and columns correspond to defined brain regions, and the entries reflect connection strengths between those regions \cite{taoudi2022predicting}. Since connectomes may differ between preterm and term births, we explore two complementary modeling perspectives: one that operates directly on the matrix representation of $A_i$, and another that interprets it as a graph structure.



\subsection{Matrix-Based Approach}

Each structural connectome is encoded as a symmetric adjacency matrix $A_i \in \mathbb{R}^{n \times n}$, where the entry $A_i[u, v]$ gives the strength of the connection between region $u$ and region $v$. Since $A_i$ is symmetric and its diagonal entries carry no information (they represent self-connections), we extract a feature vector by listing all entries above the main diagonal, where $u < v$:

\begin{equation}
    \mathbf{x}_i = \bigl(A_i[1,2],\, A_i[1,3],\, \dots,\, A_i[n-1,n]\bigr).
\end{equation}
This vector contains exactly one entry for each unordered pair of regions, capturing all unique connection strengths in the brain network. The resulting feature vector has length $d = \frac{n(n-1)}{2}$.

\subsection{Graph-Based Approach}

Alternatively, we interpret each structural connectome as a weighted, undirected graph $G_i = \bigl(V,\,E_i,\,X_v^{(i)},\,X_e^{(i)}\bigr)$, constructed from its connectivity matrix $A_i \in \mathbb{R}^{n \times n}$:

\begin{itemize}
    \item $V$ is the fixed set of $n$ nodes, each corresponding to a brain region.
    \item $E_i \subseteq V \times V$ is the subject-specific edge set, defined by nonzero connections in $A_i$.
    \item $X_v^{(i)} \in \mathbb{R}^{n \times f}$ is the feature matrix for the i-th node.
    \item $X_e^{(i)} \in \mathbb{R}^{|E_i| \times 1}$ is the edge feature matrix. Each edge $(v_j, v_k) \in E_i$ carries a connectivity strength equal to $A_i[j,k]$.
\end{itemize}

This formulation transforms each structural connectome $i$ into a graph $G_i$, enabling the use of graph-based learning models to perform binary classification. 

\subsection{Feature Augmentation and Preprocessing strategies}

To improve model performance, we explore both feature augmentation and preprocessing strategies. Table \ref{tab:augmentation_strategies} summarizes the techniques applicable to both matrix-based and graph-based formulations. Figure \ref{fig:brainconnectivity_combined} shows brain connectivity visualizations using spatial coordinates from Table \ref{tab:augmentation_strategies}.

\subsubsection{Spatial coordinates}
In order to obtain spatial coordinates we calculate node centroids based on each brain region in the infant brain atlas \cite{shi2011infant, schuh2018unbiased} which was used in the initial processing of DTI scans. We first identify unique non-zero integer labels within the atlas NIfTI image. For each brain region (1-90) we identify all voxels belonging to that region, and calculate the centroid of each giving us a mean X,Y,Z node coordinate per region. Applying the affine transformation matrix from the NIfTI header, we are able to obtain the real world spatial coordinates (in mm) from the voxel space centroids. These coordinates are approximate central locations for the corresponding nodes in the already processed connectivity matrices. This allows us to take information such as the organisation of nodes relative to one another. The resulting centroids are shown in Figure \ref{fig:brainconnectivity_combined}.

\begin{table}[htbp]
\centering
\vspace{0.125cm}
\caption{Feature augmentations and preprocessing strategies.}

\setlength{\tabcolsep}{4pt}  
\begin{tabular}{@{}p{2.6cm}p{5.6cm}@{}}
\toprule
\textbf{Feature} & \textbf{Description} \\
\midrule
Spatial Coordinates & Adds $(x, y, z)$ atlas-based centroids as node features, providing spatial information. \\
Node Degree & Adds the number of edges per node: $d_v = \sum_u \mathbb{I}[A[v,u] > 0]$, encoding local connectivity. \\
Thresholding & Removes weak edges below a threshold $\epsilon$: $A[v,u] = 0$ if $A[v,u] < \epsilon$. \\
Undersampling & Balances class distribution by randomly removing samples from the majority class. \\
\bottomrule
\end{tabular}
\label{tab:augmentation_strategies}
\end{table}







% tried with balanced/undersampling and also with stratified sampling. Also did k-folds analysis
% tried a variety of models : log regression;svm;decision tree;random forest; naiive bayes;knn;mlp

%explain coordinates because topological information is useful too.
%tried GCN; GAT/eGAT; GAT with coordinates; 

%We take coordinate information for each node as the centrality measure of each brain region in the dHCP atlas \cite{TAOUDIBENCHEKROUN2022119319}. This gives an approximation of the location of each node relative to one another within the network. 
% without the coordinate info the nodes are labelled 1-90 (based on the numerical order in the atlas) and the edges are connections between nodes, with information about how strong these connections are.



\subsection{Fuzzy logic}\label{fuzzylogic}

In the standard formulation, the gestational age (GA) label is defined as $y_i = \text{preterm}$ if $\mathrm{GA}_i < \tau$ and $y_i = \text{term}$ if $\mathrm{GA_i} \geq \tau$, with the cutoff set at $\tau = 37$ weeks \cite{quinn2016preterm}. However, GA is recorded in whole weeks, and biological maturation is continuous, so infants at 36$^{+6 days}$ and 37$^{+0 days}$ weeks often exhibit nearly identical connectomes \cite{jukic2013length}. Recent studies have shown that functional brain development evolves gradually around this gestational threshold, without a sharp boundary \citep{hou2024sensitivity}. As a result, models trained on these hard labels tend to struggle for subjects with $\mathrm{GA_i}$ near $\tau$, where small dating errors introduce label noise and the decision boundary becomes arbitrary. 

To address this, we define a soft target:

\begin{equation}
y_i^{\mathrm{soft}}
= \sigma\!\left(\frac{\mathrm{GA}_i - \tau}{T}\right)
= \frac{1}{1 + \exp\left(-\frac{\mathrm{GA}_i - \tau}{T}\right)},
\end{equation}

where $T$ is a temperature parameter that controls the smoothness of the transition, as shown in Figure \ref{fig:fuzzy_function}.

\begin{figure}[H]
    \centering
    \includegraphics[width=0.75\columnwidth]{ecai-template/figures/fuzzyfunction.png}
    \caption{Sigmoidal mapping from gestational age to soft target $y_i^{\mathrm{soft}}$. This illustrates the variables transition point ($\tau$) and steepness ($T$) and how they interact with GA. From this interaction we determine probability, and this serves as the fuzzy label.\vspace{0.9cm}}
    \label{fig:fuzzy_function}
\end{figure}


\begin{figure*}[h]
    \centering
    \begin{subfigure}[t]{0.45\textwidth} 
        \includegraphics[width=\textwidth]{ecai-template/figures/fig_fuzzy1.png}
        \label{subfig:fuzzy1}
    \end{subfigure}
    \hfill
    \begin{subfigure}[t]{0.45\textwidth}
        \includegraphics[width=\textwidth]{ecai-template/figures/fig_fuzzy2.png}
        \label{subfig:fuzzy2}
    \end{subfigure}
    \caption{The left-hand image shows traditional binary labels, where all subjects are classified as 0 (preterm) or 1 (term) based on the 37-week threshold. The right-hand image shows the same subjects after applying fuzzy logic, with a smoothed transition around the threshold. \vspace{0.5cm}}
    \label{fig:fuzzy_combined}
\end{figure*}
Training with these soft targets smooths the decision boundary around 37 weeks, which may reduce the impact of GA-recording errors. This approach encourages the model to learn graded changes in connectivity rather than an abrupt jump, and it aims to improve robustness in the late preterm window (35–37 weeks), where clinical and connectomic differences lie on a continuum. A comparison of the fuzzy and traditional boundaries is shown in Figure \ref{fig:fuzzy_combined}.


\begin{comment}
\begin{figure}[h]
    \centering
    \begin{subfigure}[t]{0.49\linewidth} 
        \includegraphics[width=\linewidth]{ecai-template/figures/fig_fuzzy1.png}
        %\caption*{}  % opcional: título vacío
        \label{subfig:fuzzy1}
    \end{subfigure}
    \hfill
    \begin{subfigure}[t]{0.49\linewidth}
        \includegraphics[width=\linewidth]{ecai-template/figures/fig_fuzzy2.png}
        %\caption*{}  % opcional: título vacío
        \label{subfig:fuzzy2}
    \end{subfigure}
    \caption{The left-hand image shows traditional binary labels, where all subjects are classified as 0 (preterm) or 1 (term) based on the 37-week threshold. The right-hand image shows the same subjects after applying fuzzy logic, with a smoothed transition around the threshold.}
    \label{fig:fuzzy_combined}
\end{figure}
\end{comment}





% Length of human pregnancy and contributors to its natural variation - Jukic et al





%maybe some plots of different T values? #temperature
%The group of infants sometimes refered to as 'late preterm infants' are often the least understood or researched, but may also be succeptable to differences from their full term counterparts born only days or weeks later.
%the 37 week mark. and also include info on the measures like spectral radius 



\subsection{Explainability}

%DB

While high classification performance demonstrates that structural connectomes carry discriminative information, understanding why a model makes its decisions is important for trust and biological insight. In brain networks, explainability reveals which connections or regions drive the prediction of preterm versus term status, guiding neuroscientific hypotheses and potential biomarkers.

We employ SHapley Additive exPlanations (SHAP) \cite{gilpin2018explaining} to decompose the prediction $f(\mathbf{x})$ as:
\begin{equation}
  f(\mathbf{x}) = \phi_{0} + \sum_{j=1}^{d} \phi_{j},
\end{equation}
where $\phi_{0}$ is the baseline output and $\phi_{j}$ is the SHAP value of feature $j$.  

After computing the SHAP matrix $\Phi \in \mathbb{R}^{N \times d}$ for $N$ subjects and $d$ features (edges and global covariates), we identify the most important edges and nodes as follows:

\begin{enumerate}
  \item \textbf{Edge importance matrix.}  
    We compute the mean absolute SHAP value for each edge feature:
    \begin{equation} \label{eq1:edge}
  M_{uv} = \frac{1}{N} \sum_{i=1}^{N} \bigl| \phi_{j(i,u,v)}^{(i)} \bigr|,
\end{equation}
where $j(i,u,v)$ maps the pair of regions $(u,v)$ to the corresponding index in $\mathbf{x}$. We then symmetrize the matrix by setting $M_{uv} = M_{vu}$.

  \item \textbf{Node-level aggregation.}  
    For each node $v$, we define:
    \begin{equation} \label{eq1:node}
  I_{v} = \frac{1}{n-1} \sum_{u \neq v} M_{uv},
\end{equation}
    which captures the average importance of all edges incident on $v$.
\end{enumerate}

For SHAP importance threshold $\delta$, edges with $M_{uv} > \delta$ and nodes with top $I_v$ values identify the most influential connections and regions driving the model's decisions. This explainability pipeline highlights the features that most strongly inform the classifier \cite{lundberg2017unified}.


% \subsubsection{Diverse Counterfactual Explanations (DiCE)}
% While SHAP considers global information about influence in model decisions, DiCE considers more on an individual case basis. The aim is to identify which features of an individual case could be altered to reverse the decision of the classifier.

\begin{comment}
\cite{molnar2020interpretable,gilpin2018explaining}.  
\cite{shapley1953value}.  
\cite{lundberg2017unified}.  
\cite{lundberg2017unified,ribeiro2016should}.
\end{comment}

%%%%%%%%%%%%%%%%%%%%%%%%%%%%%%%%%%%%%%%%%%%%%%%%%%%%%%%%%%%%%%%%%%%%%%%%


\section{Experiments}

%read the comment a decide below and decide (Alberto D.)
\subsection{Dataset}
We use data from the Developing Human Connectome Project 2nd release \cite{hughes2017dedicated} which was processed in \cite{taoudi2022predicting}, and made available on their GitHub \footnote{\url{https://github.com/CoDe-Neuro/Predicting-age-and-clinical-risk-from-the-neonatal-connectome}}. It comprises structural brain data from 524 infants, shortly after birth, acquired via DTI. Connectivity is represented as symmetric adjacency matrices between 90 cortical and subcortical regions. For each subject, we also have the following metadata: gestational age at birth (GA), postmenstrual age at scan (PMA), sex, session ID, and subject ID. We derive the preterm/term label using GA$\leq$37 weeks as preterm and GA$>$37 weeks as term.
Table \ref{tab:raw_dataset} summarizes all variables, including their array shape and type.


\begin{table}[h]
    \centering
    \caption{Description of dataset variables.}
    \footnotesize
    \begin{tabular}{lll}
        \toprule
        \textbf{Variable} & \textbf{Shape} & \textbf{Description} \\
        \midrule
        SCmu & (90, 90, 524) & Connectivity matrices for 524 subjects \\
        ga   & (524, 1)     & Gestational age at birth (weeks) \\
        mu   & (1, 524)     & Mean edge weight\\
        pma  & (524, 1)     & Postmenstrual age at scan (weeks) \\
        ses  & (1, 524)     & Session identifiers (string) \\
        sex  & (524, 1)     & Sex (0 = female, 1 = male) \\
        sub  & (1, 524)     & Subject identifiers (string) \\
        \bottomrule
    \end{tabular}
    \label{tab:raw_dataset}
\end{table}


\subsection{Design}

We apply the proposed methodology to the dataset described above. We use four baseline ML models: Logistic Regression (LR), SVM, Multi‑Layer Perceptron (MLP), and Random Forest (RF) \cite{mohamed2023hybrid}. We also evaluate two GNN architectures: Graph Convolutional Network (GCN) and Graph Attention Network (GAT) \cite{zhang2019graph, velivckovic2017graph}. Each execution uses 5‑fold cross-validation (CV), with hyperparameters selected as shown in Table \ref{tab:augmentation_strategies}. We  use stratified sampling and use a 20\% / 80\% training/test split. Additionally, we integrate fuzzy logic into the best‑performing models from both the ML and GNN approaches. We use accuracy, precision, recall and F1‑score as the evaluation metrics for all experiments.
In addition, we calculate weighted averages, in order to account for unbalanced classes \cite{sokolova2009systematic}:

\begin{equation}
  \text{Metric}_{W} = \frac{(\text{Metric}_{0} \times \text{Class size}_{0}) + (\text{Metric}_{1} \times \text{Class size}_{1})}{\text{Class size}_{0} + \text{Class size}_{1}}.
\end{equation}


\subsection{Results}
In Table \ref{tab:classification_spatial} we present the ML and GNN results alongside those obtained using the spatial‑coordinate (described in Table \ref{tab:augmentation_strategies}). We focus on those atlas coordinates from Table \ref{tab:augmentation_strategies} because this strategy was the only one to improve model performance. Based on these results, we select the best‑performing algorithms from each approach (LR with spatial coordinates and GAT) and apply fuzzy‑logic labeling, as shown in Table \ref{tab:results_fuzzy}.



\begin{table*}[h]
    \centering    
    \caption{Impact of Spatial Information on ML Model Performance. We show LR, SVM, MLP, RF. We also show GNN variants: GCN, GAT, and Graph Isomorphism Network (GIN). Each model is assessed with and without the inclusion of spatial node coordinates. All values are shown from 0 to 1. Best results are in bold. Weighted averages consider unbalanced classes [21 preterm; 84 term].}
    \vspace{0.1cm}

    \begin{tabular}{ll|ccc|cccc}
        \hline
        \textbf{Model} & \textbf{Class} & \textbf{Precision} & \textbf{Recall} & \textbf{F1-score}
            & \multicolumn{4}{c}{\textbf{Overall (weighted average)}} \\
        & & & & & \textbf{Accuracy} & \textbf{Precision} & \textbf{Recall} & \textbf{F1-score} \\
        \hline
        \multirow{2}{*}{LR} 
            & Preterm & 0.70& 0.76& 0.73& \multirow{2}{*}{0.89}& \multirow{2}{*}{0.89}& \multirow{2}{*}{0.89}& \multirow{2}{*}{0.89}\\
            & Term    & 0.94& 0.92& 0.93&  &  &  &  \\
        \hline
        \multirow{2}{*}{LR + Spatial} 
            & Preterm    & 0.85& 0.81& \textbf{0.83}& \multirow{2}{*}{\textbf{0.93}}& \multirow{2}{*}{\textbf{0.93}}& \multirow{2}{*}{\textbf{0.93}}& \multirow{2}{*}{\textbf{0.93}}\\
            & Term       & 0.95& 0.96& \textbf{0.96}&  &  &  &  \\
        \hline
        \multirow{2}{*}{SVM} 
            & Preterm & 0.61& 0.90& 0.73& \multirow{2}{*}{0.87}& \multirow{2}{*}{0.90}& \multirow{2}{*}{0.87}& \multirow{2}{*}{0.88}\\
            & Term    & \textbf{0.97}& 0.86& 0.91&  &  &  &  \\
        \hline
        \multirow{2}{*}{SVM + Spatial} 
            & Preterm & 0.20& \textbf{1.00}& 0.33& \multirow{2}{*}{0.20}& \multirow{2}{*}{0.04}& \multirow{2}{*}{0.20}& \multirow{2}{*}{0.07}\\
            & Term    & 0.00& 0.00& 0.00&  &  &  &  \\
        \hline
        \multirow{2}{*}{MLP} 
            & Preterm & 0.60& 0.14& 0.23& \multirow{2}{*}{0.81}& \multirow{2}{*}{0.78}& \multirow{2}{*}{0.81}& \multirow{2}{*}{0.76}\\
            & Term    & 0.82& 0.98& 0.89&  &  &  &  \\
        \hline
        \multirow{2}{*}{MLP + Spatial} 
            & Preterm & 0.00& 0.00& 0.00& \multirow{2}{*}{0.80}& \multirow{2}{*}{0.64}& \multirow{2}{*}{0.80}& \multirow{2}{*}{0.71}\\
            & Term    & 0.80& \textbf{1.00}& 0.89&  &  &  &  \\
        \hline
        \multirow{2}{*}{RF} 
            & Preterm & \textbf{1.00}& 0.19& 0.32& \multirow{2}{*}{0.84}& \multirow{2}{*}{0.87}& \multirow{2}{*}{0.84}& \multirow{2}{*}{0.79}\\
            & Term    & 0.83& \textbf{1.00}& 0.91&  &  &  &  \\
        \hline
        \multirow{2}{*}{RF + Spatial} 
            & Preterm & \textbf{1.00}& 0.10& 0.17& \multirow{2}{*}{0.82}& \multirow{2}{*}{0.85}& \multirow{2}{*}{0.82}& \multirow{2}{*}{0.75}\\
            & Term    & 0.82& \textbf{1.00}& 0.90&  &  &  &  \\
        \hline\hline
        \multirow{2}{*}{GCN}
            & Preterm & \textbf{0.91}& 0.48& 0.63& \multirow{2}{*}{0.89}& \multirow{2}{*}{0.89}& \multirow{2}{*}{0.89}& \multirow{2}{*}{0.87}\\
            & Term    & 0.88& 0.99& 0.93&                        &                        &                        &                        \\
        \hline
        \multirow{2}{*}{GCN + Spatial}
            & Preterm & 0.00& 0.00& 0.00& \multirow{2}{*}{0.80}& \multirow{2}{*}{0.64}& \multirow{2}{*}{0.80}& \multirow{2}{*}{0.71}\\
            & Term    & 0.80& \textbf{1.00}& 0.89&                        &                        &                        &                        \\
        \hline
        \multirow{2}{*}{GAT}
            & Preterm & 0.79& \textbf{0.71}& \textbf{0.75}& \multirow{2}{*}{\textbf{0.90}}& \multirow{2}{*}{\textbf{0.90}}& \multirow{2}{*}{\textbf{0.90}}& \multirow{2}{*}{\textbf{0.90}}\\
            & Term    & \textbf{0.93}& 0.95& \textbf{0.94}&                        &                        &                        &                        \\
        \hline
        \multirow{2}{*}{GAT + Spatial}
            & Preterm & 0.85& 0.52& 0.65& \multirow{2}{*}{0.89}& \multirow{2}{*}{0.88}& \multirow{2}{*}{0.89}& \multirow{2}{*}{0.87}\\
            & Term    & 0.89& \textbf{0.98}& 0.93&                        &                        &                        &                        \\
        \hline
        \multirow{2}{*}{GIN}
            & Preterm & 0.00 & 0.00 & 0.00 &
            \multirow{2}{*}{0.80}&
            \multirow{2}{*}{0.64}&
            \multirow{2}{*}{0.80}&
            \multirow{2}{*}{0.71}\\
            & Term & 0.80 & \textbf{1.00} & 0.89 & 
            &
            &
            &
            \\
            \hline
        \multirow{2}{*}{GIN + Spatial}
            & Preterm & 0.00 & 0.00 & 0.00 &
            \multirow{2}{*}{0.80}&
            \multirow{2}{*}{0.64}&
            \multirow{2}{*}{0.80}&
            \multirow{2}{*}{0.71}\\
            & Term & 0.80 & \textbf{1.00} & 0.89 & 
            &
            &
            &
            \\
            \hline
    \end{tabular}
    \label{tab:classification_spatial}
\end{table*}




\begin{table*}[h]
    \centering
    \caption{Performance metrics with fuzzy logic. We select the two best performing models from Table above: LR with Spatial and GAT. Values are shown from 0 to 1 with four decimal places.}
    \vspace{0.1cm}
    \begin{tabular}{ll|ccc|cccc}
        \hline
        \textbf{Model} & \textbf{Class} & \textbf{Precision} & \textbf{Recall} & \textbf{F1-score}
            & \multicolumn{4}{c}{\textbf{Overall (weighted average)}} \\
        & & & & & \textbf{Accuracy} & \textbf{Precision} & \textbf{Recall} & \textbf{F1-score} \\
        \hline
        
        \multirow{2}{*}{LR + Spatial + fuzzy} 
            & Preterm    & \textbf{0.95}& \textbf{0.86}& \textbf{0.90}& \multirow{2}{*}{\textbf{0.96}}& \multirow{2}{*}{\textbf{0.96}}& \multirow{2}{*}{\textbf{0.96}}& \multirow{2}{*}{\textbf{0.96}}\\
            & Term       & \textbf{0.97}& \textbf{0.99}& \textbf{0.98}&  &  &  &  \\
        \hline
        \multirow{2}{*}{GAT + fuzzy} 
            & Preterm    & 0.92& 0.52& 0.67& \multirow{2}{*}{0.90}& \multirow{2}{*}{0.90}& \multirow{2}{*}{0.90}& \multirow{2}{*}{0.88}\\
            & Term       & 0.89& \textbf{0.99}& 0.94&  &  &  &  \\
        \hline
    \end{tabular}
    \label{tab:results_fuzzy}
    \vspace{0.05cm}
\end{table*}







% \begin{table}[h!]
%     \caption{ML models baselines: Random Forest (RF); Support Vector Machine (SVM); MultiLayer Perceptron (MLP); Logistic Regression (LR). For all models, support is 21 (preterm) and 84 (term)}
%     \centering
%     \begin{tabularx}{\columnwidth}{lXXXX}
        
%         \hline
%         % \multirow{2}{*}{Metric} & \multirow{2}{*}{Logistic Regression} & \multirow{2}{*}{SVM} & \multirow{2}{*}{MLP} & \multirow{2}{*}{Random Forest} \\
%         % & & & & \\ 
%         Metric                 & Logistic Regression & SVM       & MLP       & Random Forest \\
%         \hline
%         % \textbf{Training CV Accuracy} & \textbf{90.69 \%}  & 84.49 \%  & 79.71 \%  & 84.72 \%      \\
%         \textbf{Accuracy}             & 88.57 \%    & 86.67 \%   & 80.95 \%  & 83.14 \%      \\
%         \hline
%         \textbf{Precision (Pre-term)} & 69.57 \%    & 61.29 \%   & 60.00 \%  & 100.00 \%     \\
%         \textbf{Recall (Pre-term)}    & 76.19 \%    & 90.48 \%   & 14.29 \%  & 19.05 \%      \\
%         \textbf{F1-Score (Pre-term)}  & 0.73        & 0.73       & 0.23      & 0.32          \\
%         \hline
%         \textbf{Precision (Term)}     & 93.90 \%    & 97.30 \%   & 82.00 \%  & 83.17 \%      \\
%         \textbf{Recall (Term)}        & 91.67 \%    & 85.71 \%   & 97.62 \%  & 100.00 \%     \\
%         \textbf{F1-Score (Term)}      & 0.93        & 0.91       & 0.89      & 0.91          \\
%         \hline
%         \textbf{AUC}                & 0.89          & 0.90       & 0.78      & 0.94          \\
        
%     \end{tabularx}

%     \label{tab:model_comparison}
%     \footnotesize{}
% \end{table}



% \begin{table}[h!]
%     \caption{ML models with coordinates: Random Forest (RF); Support Vector Machine (SVM); MultiLayer Perceptron (MLP); Logistic Regression (LR). For all models, support is 21 (preterm) and 84 (term)}
%     \centering
%     \begin{tabularx}{\columnwidth}{lXXXX}
        
%         \hline
%         Metric                 & Logistic Regression & SVM       & MLP       & Random Forest \\
%         \hline
%         % \textbf{Training CV Accuracy}       & \textbf{88.07 \%}  & 79.71 \%  & 79.71 \%  & 85.20 \%      \\
%         \textbf{Accuracy}             & 93.33 \%    & 20.00 & 80.00 \% & 81.90 \% \\
%         \hline
%         \textbf{Precision (Pre-term)} & 85.00 \%    & 20.00 \%   & 00.00 \%  & 100.00 \%     \\
%         \textbf{Recall (Pre-term)}    & 80.95 \%    & 100.00 \%   & 00.00 \%  & 9.05 \%      \\
%         \textbf{F1-Score (Pre-term)}  & 0.83        & 0.33       & 0.00      & 0.17          \\
%         \hline
%         \textbf{Precision (Term)}     & 95.29 \%    & 00.00 \%   & 80.00 \%  & 81.55 \%      \\
%         \textbf{Recall (Term)}        & 96.43 \%    & 00.00 \%   & 100.00 \%  & 100.00 \%     \\
%         \textbf{F1-Score (Term)}      & 0.96        & 0.00       & 0.89      & 0.90          \\
%         \hline
%         \textbf{AUC}                & 0.91          & 0.06       & 0.50      & 0.90          \\
        
%     \end{tabularx}

%     \label{tab:model_comparison}
%     \footnotesize{}
% \end{table}


\begin{comment}
\begin{table}[h!]
    \centering
    \renewcommand{\arraystretch}{1.3}
    \caption{Impact of Spatial Information on ML Model Performance. First number for each metric is the baseline model, and the second value is the model with addition of spatial coordinate information. We indicate the change in brackets.}
    \vspace{0.1cm}

    \begin{tabularx}{\columnwidth}{X
                                     >{\hsize=0.8\hsize\arraybackslash}X
                                     >{\hsize=1.1\hsize\arraybackslash}X 
                                     >{\hsize=1.1\hsize\arraybackslash}X 
                                     >{\hsize=1.0\hsize\arraybackslash}X} 
        \hline
        \textbf{Metric} & \textbf{LR} & \textbf{SVM} & \textbf{MLP} & \textbf{RF} \\
        \hline
        Accuracy & \textbf{88.57\%} & 86.67\% & 80.95\% & 83.14\% \\
                         & \textbf{93.33\%} & 20.00\% & 80.00\% & 81.90\% \\
                         & \textbf{(+4.76\%)} & (-66.67\%) & (-0.95\%) & (-1.24\%) \\
        \hdashline
        AUC & 0.89 & 0.90 & 0.78 & \textbf{0.94 }\\
                     &\textbf{ 0.90 }& 0.06 & 0.50 & 0.90 \\
                     & \textbf{(+0.01) }& (-0.84) & (-0.28) & (-0.04) \\
        \hline
       \textbf{Preterm} & & & & \\
        \hline
        Precision & 69.57\% & 61.29\% & 60.00\% & \textbf{100.00\%} \\
                               & 85.00\% & 20.00\% & 00.00\% & \textbf{100.00\%} \\
                               & \textbf{(+15.43\%)} & (-41.29\%) & (-60.00\%) & (+0.00\%) \\
        \hdashline
        Recall & 76.19\% & \textbf{90.48\%} & 14.29\% & 19.05\% \\
                              & 80.95\% & \textbf{100.00\%} & 00.00\% & 9.05\% \\
                              & (+4.76\%) & \textbf{(+9.52\%)} & (-14.29\%) & (-10.00\%) \\
        \hdashline
        F1-Score &\textbf{ 0.73} & \textbf{0.73} & 0.23 & 0.32 \\
                                 & \textbf{0.83} & 0.33 & 0.00 & 0.17 \\
                                 & \textbf{(+0.10)} & (-0.40) & (-0.23) & (-0.15) \\
        \hline
        \textbf{Term} & & & & \\
        \hline
        Precision & 93.90\% &\textbf{ 97.30\%} & 82.00\% & 83.17\% \\
                              & \textbf{95.29\%} & 00.00\% & 80.00\% & 81.55\% \\
                              & \textbf{(+1.39\%)} & (-97.30\%) & (-2.00\%) & (-1.62\%) \\
        \hdashline
        Recall & 91.67\% & 85.71\% & 97.62\% & \textbf{100.00\%} \\
                           & 96.43\% & 00.00\% & \textbf{100.00\%} & \textbf{100.00\%} \\
                           & \textbf{(+4.76\%)} & (-85.71\%) & (+2.38\%) & (+0.00\%) \\
        \hdashline
        F1-Score & \textbf{0.93} & 0.91 & 0.89 & 0.91 \\
                               &\textbf{ 0.96} & 0.00 & 0.89 & 0.90 \\
                               & \textbf{(+0.03}) & (-0.91) & (+0.00) & (-0.01) \\
        \hline
    \end{tabularx}
    \label{tab:spatial_impact}
\end{table}
\end{comment}


% \begin{table}[h]
%     \centering
%     \setlength{\tabcolsep}{4pt}
%     \begin{tabular}{@{}lccc@{}}
%         \toprule
%         & \textbf{accuracy} & \textbf{f1 (weighted)} & \textbf{AUC}\\
%         \midrule
%         Overall & \makecell{92.37\% \\ $\pm$ 2.32 \%} & \makecell{95.12\% \\ $\pm$ 1.56\%} & \makecell{0.973 \\ $\pm$ 0.016} \\
%         \toprule
%         & \textbf{precision} & \textbf{recall} & \textbf{Support}\\
%         \midrule
%         Preterm (C0) & \makecell{78.79\% \\ $\pm$ 7.34\%} & \makecell{87.06\% \\ $\pm$ 6.86\%} & \makecell{17.00 \\ $\pm$ 0.00} \\
%         Term (C1) & \makecell{96.65\% \\ $\pm$ 1.68\%} & \makecell{93.72\% \\ $\pm$ 3.03\%} & \makecell{66.80 \\ $\pm$ 0.40} \\
%         \bottomrule
%         \newline
%     \end{tabular}
%     \caption{Results of 5 fold cross-validation for logistic regression model with fuzzy boundary}
%     \label{tab:my_label}
% \end{table}




% \begin{table}[h]
%     \centering
%     \setlength{\tabcolsep}{4pt}
%     \begin{tabular}{@{}lcccccc@{}}
%     \toprule
%       & \textbf{accuracy} & \textbf{f1 (weighted)} & \textbf{AUC}\\
%     \midrule
%     Overall & \makecell{90.69\% \\ $\pm$ 2.07 \%} & \makecell{94.16\% \\ $\pm$ 1.32\%} & \makecell{0.957 \\ $\pm$ 0.016} \\
%     \toprule
%     {} & \textbf{precision} & \textbf{recall} & \textbf{Support}\\
%     \midrule
%     Preterm (C0) & \makecell{77.74\% \\ $\pm$ 6.24\%} & \makecell{76.47\% \\ $\pm$ 5.26\%} & \makecell{17.00 \\ $\pm$ 0.00} \\
%     Term (C1) & \makecell{94.04\% \\ $\pm$ 1.29\%} & \makecell{94.31\% \\ $\pm$ 1.98\%} & \makecell{66.8 \\ $\pm$ 0.40} \\
%     \bottomrule
%     \newline
%     \end{tabular}
%     \caption{Results of 5 fold cross-validation for logistic regression model }
%     \label{tab:my_label}
% \end{table}

% \begin{table}[htbp]
%     \centering
%     \setlength{\tabcolsep}{4pt}
%     \begin{tabular}{@{}lcccccc@{}}
%     \toprule
%     {} & \textbf{accuracy} & \textbf{f1 (weighted)} & \textbf{AUC}\\
%     \midrule
%      Overall & 88.57\% & 92.77\% & 0.892\\
%      \toprule
%     {} & \textbf{precision} & \textbf{recall} & \textbf{f1}\\
%     \midrule
%      Preterm (C0) & 69.57\% & 76.19\% & 72.73 \\
%      Term (C1) & 93.9\% & 91.67\% & 92.77\% \\
%      \bottomrule
%      \newline

%     \end{tabular}
%     \caption{Baseline Results on unseen test set}
%     \label{tab:my_label}
% \end{table}]







\begin{comment}
\begin{table}[h!]
    \centering
    \renewcommand{\arraystretch}{1.3}
    \caption{Logistic regression model and GAT model with fuzzy boundaries. \textit{NB/ LR model is marginally better with fuzzy and no spatial, while GAT requires spatial}}
    \vspace{0.1cm}
    % \begin{tabularx}{\columnwidth}{lXXXX}
    %                                  % % >{\hsize=0.8\hsize\centering\arraybackslash}X
    %                                  % >{\hsize=2.0\hsize\centering\arraybackslash}X 
    %                                  % % >{\hsize=1.8\hsize\centering\arraybackslash}X
    %                                  % % >{\hsize=1.8\hsize\centering\arraybackslash}X}
    %                                  % XX}
    %     \hline >{\hsize=0.5\hsize\arraybackslash}>{\hsize=1.0\hsize\arraybackslash}X
    \begin{tabularx}{\columnwidth}{XXX}
        \hline
        \textbf{Metric} & \textbf{Logistic Regression} & \textbf{GAT} \\
        \hline
        Accuracy & \textbf{96.19\%} & \% \\
        \hdashline
        AUC & 0.96 & 0. \\
        \hline
       \textbf{Preterm} & & \\
        \hline
        Precision & 94.74\% & \% \\
        \hdashline
        Recall & 85.71\% & \%\\
        \hdashline
        F1-Score &\textbf{0.90} & \\
        \hline
        \textbf{Term} & & \\
        \hline
        Precision & 96.51\% &\%\\
        \hdashline
        Recall & 98.81\% & \% \\
        \hdashline
        F1-Score & 0.98 &  \\
        \hline
    \end{tabularx}
    \label{tab:fuzzy_impact}
\end{table}

% \begin{table}[h]
%     \centering
%     \setlength{\tabcolsep}{4pt}
%     \begin{tabular}{@{}lccc@{}}
%         \toprule
%         & \textbf{accuracy} & \textbf{f1 (weighted)} & \textbf{AUC}\\
%         \midrule
%         Overall & \makecell{92.37\% \\ $\pm$ 2.32 \%} & \makecell{95.12\% \\ $\pm$ 1.56\%} & \makecell{0.973 \\ $\pm$ 0.016} \\
%         \toprule
%         & \textbf{precision} & \textbf{recall} & \textbf{Support}\\
%         \midrule
%         Preterm (C0) & \makecell{78.79\% \\ $\pm$ 7.34\%} & \makecell{87.06\% \\ $\pm$ 6.86\%} & \makecell{17.00 \\ $\pm$ 0.00} \\
%         Term (C1) & \makecell{96.65\% \\ $\pm$ 1.68\%} & \makecell{93.72\% \\ $\pm$ 3.03\%} & \makecell{66.80 \\ $\pm$ 0.40} \\
%         \bottomrule
%         \newline
%     \end{tabular}
%     \caption{Results of 5 fold cross-validation for logistic regression model with fuzzy boundary}
%     \label{tab:my_label}
% \end{table}
\end{comment}

% % 

% \begin{table}[h]
%     \centering
%     \setlength{\tabcolsep}{4pt}
%     \begin{tabular}{@{}lcccccc@{}}
%     \toprule
%     {} & \textbf{accuracy} & \textbf{f1 (weighted)} & \textbf{AUC}\\
%     \midrule
%      Overall & 96.19\% & 97.65\% & 0.963\\
%      \toprule
%     {} & \textbf{precision} & \textbf{recall} & \textbf{f1}\\
%     \midrule
%      Preterm (C0) & 94.74\% & 85.71\% & 90.00\% \\
%      Term (C1) & 96.51\% & 98.81\% & 97.65\% \\
%     \end{tabular}
%     \caption{Results on unseen test set using fuzzy boundary}
%     \label{tab:my_label}
% \end{table}


\subsection{Explainability}

We use SHAP Equations \ref{eq1:edge} and \ref{eq1:node} to compute attribution scores, deriving edge‑importance matrices and aggregating node‑level importances to interpret each model prediction. Figure \ref{fig:combined_importance} contrasts feature importance in the LR and GAT models through both heatmaps and Atlas-based network plots. The top row displays SHAP edge‐importance heatmaps, highlighting which connections most drive each model’s predictions. The bottom row renders nodes at their Atlas coordinates, sized and colored by SHAP score to reveal the most (red) and least (green) influential regions (node indices correspond to Table \ref{tab:regions_best_worst_overallheader}). Figure \ref{fig:connectionsshap} presents the SHAP summary plot indicating influence of individual connections on the LR model outcome. Negative SHAP values indicate the model is pushed towards class 0 (preterm) while positive values classify towards class 1 (term). Red indicates high feature values while blue indicate low. Finally, Table \ref{tab:regions_best_worst_overallheader} presents the top 10 and bottom 10 nodes ranked by SHAP importance for both the LR and GAT models. Regions shown in bold indicate agreement between the two models. LR and GAT concur on 8 of the 10 least important nodes, whereas they align on 4 of the 10 most important regions.

% \begin{table*}[h]
%     \centering
%     \caption{Most and least important nodes, as indicated by SHAP analysis. Comparison from both LR and GAT models. We highlight in bold the regions which are consistent across each model. From this we see that both LR and GAT identify 8 out of 10 same least important nodes in the decisions. For most important regions, there is consistency in 4/10 regions. }
%     \scriptsize
%     \renewcommand{\arraystretch}{1.25}
%     \begin{tabular}{@{}ll|ll@{}}
%     % \begin{tabular}{ll|ll}
%     % \begin{tabularx}{\columnwidth}{m{0.20\columnwidth} m{0.20\columnwidth}| m{0.24\columnwidth} m{0.24\columnwidth}}
%     \toprule
%         \multicolumn{2}{c}{\textbf{Top 10}} & \multicolumn{2}{c}{\textbf{Bottom 10}} \\
%         \cmidrule(lr){1-2} \cmidrule(lr){3-4}
%         \textbf{LR} & \textbf{GAT} & \textbf{LR} & \textbf{GAT} \\        \midrule
%         \textbf{77 - Thalamus R}        & 89  - Temporal Inf R & \textbf{53 - Occipital Inf R} & 49 - Occipital Sup R\\
%         30 - Cingulum Ant L &  60 - Parietal Inf L          & 20 - Olfactory L & 86 - Hippocampus L\\
%         31 - Cingulum Ant R & \textbf{6 - Frontal Mid L} & \textbf{10 - Frontal Inf Oper L} & \textbf{74 - Palladium L}\\
%         33 - Cingulum Mid R &  \textbf{29 - Insula R}          & \textbf{69 - Paracentral Lob R} & \textbf{ 36 - Frontal Inf Oper L}\\
%         35 - Cingulum Post R & 81 - Temporal Sup R  & \textbf{74 - Palladium L} & \textbf{48 - Occipital Sup L}\\
%         \textbf{29 - Insula R}                  &  \textbf{77 - Thalamus R}          & 68 - Paracentral Lob L & \textbf{ 53 - Occipital Inf R}\\
%         67 - Precuneus R          & 28 - Insula L  & \textbf{79 - Heschl } &  44 - Cuneus L\\
%         7 - Frontal Mid R   &  61 - Parietal Inf R    & \textbf{48 - Occipital Sup L} & \textbf{69 - Paracentral Lob R}\\
%         \textbf{73 - Putamen R}            &    \textbf{73 - Putamen R}  & \textbf{36 - Hippocampus L} & \textbf{ 79 - Heschl R}\\
%         \textbf{6 - Frontal Mid L}    &  70 - Caudate L            & \textbf{78 - Heschl L} & \textbf{ 78 - Heschl L}\\
%         \bottomrule
%     \end{tabular}
%     \label{tab:regions_best_worst}
% \end{table*}

\begin{table}[h]
  \centering
  \vspace{0.05cm}
  \caption{Top 10 and bottom 10 nodes by SHAP importance for both LR and GAT models, with bold indicating regions of agreement (8/10 least important; 4/10 most important. SHAP values calculated through the combined metric. These regions are displayed in Atlas coordinate space in Figure \ref{fig:combined_importance} panels (c) and (d).}
  \scriptsize
  \renewcommand{\arraystretch}{1.0}
  \begin{tabular}{@{}llcllc@{}}
    \toprule
      \multicolumn{3}{c|}{\textbf{LR}} & \multicolumn{3}{c}{\textbf{GAT}} \\
      \cmidrule(lr){1-3} \cmidrule(lr){4-6}
      \multicolumn{6}{c}{\textbf{Top 10}} \\
      % \cmidrule(x){}
      \midrule
      \textbf{\shortstack{Node \\ Index}} & \textbf{\shortstack{Node \\ Name}} & \textbf{\shortstack{SHAP \\ Value}} & \textbf{\shortstack{Node \\ Index}} & \textbf{\shortstack{Node \\ Name}} & \textbf{\shortstack{SHAP \\ Value}} \\
      % \textbf{Node  Index} & \textbf{Node Name} & \textbf{SHAP Value} & \textbf{Node  Index} & \textbf{Node  Name} & \textbf{SHAP  Value} \\
      \midrule
      \textbf{78} & \textbf{Thalamus R} & 0.22  & 90 & Temporal Inf R & 0.39\\
      31 & Cingulum Ant L & 0.21 & 61 & Parietal Inf L & 0.38\\
      32 & Cingulum Ant R & 0.20 & \textbf{7} & \textbf{Frontal Mid L} & 0.37\\
      33 & Cingulum Mid R & 0.18 & \textbf{30} & \textbf{Insula R} & 0.35\\
      36 & Cingulum Post R & 0.18 & 82 & Temporal Sup R & 0.33\\
      \textbf{30} & \textbf{Insula R} & 0.18 & \textbf{78} & \textbf{Thalamus R} & 0.33\\
      68 & Precuneus R & 0.18 & 29 & Insula L & 0.33\\
      8 & Frontal Mid R & 0.16& 62 & Parietal Inf R & 0.33\\
      \textbf{74} & \textbf{Putamen R} & 0.16 & \textbf{74} & \textbf{Putamen R} & 0.33\\
      \textbf{7} & \textbf{Frontal Mid L} & 0.16 & 71 & Caudate L & 0.31\\
      \midrule
      \multicolumn{6}{c}{\textbf{Bottom 10}}\\
      \midrule
      \textbf{54} & \textbf{Occipital Inf R} & 0.06& 50 & Occipital Sup R & 0.09\\
      21 & Olfactory L & 0.06 & 87 & Hippocampus L & 0.08\\
      \textbf{11} & \textbf{Frontal Inf Oper L} & 0.06 & \textbf{75} & \textbf{Palladium L} & 0.08\\
      \textbf{70} & \textbf{Paracentral Lob R} & 0.06 & \textbf{37} & \textbf{Frontal Inf Oper L} & 0.08\\
      \textbf{75} & \textbf{Palladium L} & 0.05 & \textbf{49} & \textbf{Occipital Sup L} & 0.07\\
      69 & Paracentral Lob L & 0.05 & \textbf{54} & \textbf{Occipital Inf R} & 0.07\\
      \textbf{80} & \textbf{Heschl } & 0.05 & 45 & Cuneus L & 0.07\\
      \textbf{49} & \textbf{Occipital Sup L} & 0.04 & \textbf{70} & \textbf{Paracentral Lob R} &0.06\\
      \textbf{37} & \textbf{Hippocampus L} & 0.04 & \textbf{80} & \textbf{Heschl R} & 0.03\\
      \textbf{79} & \textbf{Heschl L} & 0.01 & \textbf{79} & \textbf{Heschl L} & 0.03\\
    \bottomrule
  \end{tabular}
  \label{tab:regions_best_worst_overallheader}
\end{table}
% \begin{table}[h]
%   \centering
%   \caption{Top 10 and bottom 10 nodes by SHAP importance for both LR and GAT models, with bold indicating regions of agreement (8/10 least important; 4/10 most important).}
%   \scriptsize
%   \renewcommand{\arraystretch}{1.25}
%   \begin{tabular}{@{}ll@{}}
%     \toprule
%       \multicolumn{2}{c}{\textbf{Top 10}} \\
%       \cmidrule(lr){1-2}
%       \textbf{LR} & \textbf{GAT} \\
%       \midrule
%       \textbf{77 - Thalamus R} & 89 - Temporal Inf R \\
%       30 - Cingulum Ant L & 60 - Parietal Inf L \\
%       31 - Cingulum Ant R & \textbf{6 - Frontal Mid L} \\
%       33 - Cingulum Mid R & \textbf{29 - Insula R} \\
%       35 - Cingulum Post R & 81 - Temporal Sup R \\
%       \textbf{29 - Insula R} & \textbf{77 - Thalamus R} \\
%       67 - Precuneus R & 28 - Insula L \\
%       7 - Frontal Mid R & 61 - Parietal Inf R \\
%       \textbf{73 - Putamen R} & \textbf{73 - Putamen R} \\
%       \textbf{6 - Frontal Mid L} & 70 - Caudate L \\
%       \midrule
%       \multicolumn{2}{c}{\textbf{Bottom 10}} \\
%       \cmidrule(lr){1-2}
%       \textbf{LR} & \textbf{GAT} \\
%       \midrule
%       \textbf{53 - Occipital Inf R} & 49 - Occipital Sup R \\
%       20 - Olfactory L & 86 - Hippocampus L \\
%       \textbf{10 - Frontal Inf Oper L} & \textbf{74 - Palladium L} \\
%       \textbf{69 - Paracentral Lob R} & \textbf{36 - Frontal Inf Oper L} \\
%       \textbf{74 - Palladium L} & \textbf{48 - Occipital Sup L} \\
%       68 - Paracentral Lob L & \textbf{53 - Occipital Inf R} \\
%       \textbf{79 - Heschl } & 44 - Cuneus L \\
%       \textbf{48 - Occipital Sup L} & \textbf{69 - Paracentral Lob R} \\
%       \textbf{36 - Hippocampus L} & \textbf{79 - Heschl R} \\
%       \textbf{78 - Heschl L} & \textbf{78 - Heschl L} \\
%     \bottomrule
%   \end{tabular}
%   \label{tab:regions_best_worst_onecolumn}
% \end{table}
% \subsubsection{Ablations}
% KB/AD/DB
% eg - with and without the coordinate node features
%%%%%%%%%%%%%%%%%%%%%%%%%%%%%%%%%%%%%%%%%%%%%%%%%%%%

% \begin{table*}[htbp]
%     \centering
%     \caption{Most and least important nodes, as indicated by SHAP analysis. Comparison from both LR and GAT models. We highlight in bold the regions which are consistent across each model. From this we see that both LR and GAT identify 8 out of 10 same least important nodes in the decisions. For most important regions, there is consistency in 4/10 regions. }
%     \scriptsize
%     \renewcommand{\arraystretch}{1.25}
%     \begin{tabular}{@{}ll|ll@{}}
%     % \begin{tabular}{ll|ll}
%     % \begin{tabularx}{\columnwidth}{m{0.20\columnwidth} m{0.20\columnwidth}| m{0.24\columnwidth} m{0.24\columnwidth}}
%     \toprule
%         \multicolumn{2}{c}{\textbf{Top 10}} & \multicolumn{2}{c}{\textbf{Bottom 10}} \\
%         \cmidrule(lr){1-2} \cmidrule(lr){3-4}
%         \textbf{LR} & \textbf{GAT} & \textbf{LR} & \textbf{GAT} \\         \midrule
%         \textbf{77 - Thalamus R}       & 89  - Temporal Inf R & \textbf{53 - Occipital Inf R} & 49 - Occipital Sup R\\
%         30 - Cingulum Ant L &  60 - Parietal Inf L              & 20 - Olfactory L & 86 - Hippocampus L\\
%         31 - Cingulum Ant R & \textbf{6 - Frontal Mid L} & \textbf{10 - Frontal Inf Oper L} & \textbf{74 - Palladium L}\\
%         33 - Cingulum Mid R &  \textbf{29 - Insula R}          & \textbf{69 - Paracentral Lob R} & \textbf{ 36 - Frontal Inf Oper L}\\
%          35 - Cingulum Post R & 81 - Temporal Sup R     & \textbf{74 - Palladium L} & \textbf{48 - Occipital Sup L}\\
%         \textbf{29 - Insula R}               &  \textbf{77 - Thalamus R}        & 68 - Paracentral Lob L & \textbf{ 53 - Occipital Inf R}\\
%         67 - Precuneus R         & 28 - Insula L   & \textbf{79 - Heschl } &  44 - Cuneus L\\
%         7 - Frontal Mid R      &  61 - Parietal Inf R      & \textbf{48 - Occipital Sup L} & \textbf{69 - Paracentral Lob R}\\
%         \textbf{73 - Putamen R}          &   \textbf{73 - Putamen R}     & \textbf{36 - Hippocampus L} & \textbf{ 79 - Heschl R}\\
%         \textbf{6 - Frontal Mid L}    &  70 - Caudate L             & \textbf{78 - Heschl L} & \textbf{ 78 - Heschl L}\\
%         \bottomrule
%     \end{tabular}
%     \label{tab:regions_best_worst}
% \end{table*}

\begin{figure*}[h]
    \centering
    \begin{subfigure}[t]{0.45\textwidth}
        \includegraphics[width=\textwidth]{ecai-template/figures/matrixheatmap.png}
        \caption{LR (Edges)}
        \label{subfig:LRheatmap}
    \end{subfigure}%
    \hfill
    \begin{subfigure}[t]{0.45\textwidth}
        \includegraphics[width=\textwidth]{ecai-template/figures/graphheatmap.png}
        \caption{GAT (Edges)}
        \label{subfig:GATheatmap}
    \end{subfigure}
    \vfill
    \vspace{0.35cm}% Add vertical space between the rows of subfigures
    \begin{subfigure}[t]{0.45\textwidth}
        \includegraphics[width=\textwidth]{ecai-template/figures/matrixnodesSHAP.png}
        \caption{LR (Nodes)}
        \label{subfig:LRSHAP}
    \end{subfigure}%
    \hfill
    \begin{subfigure}[t]{0.45\textwidth}
        \includegraphics[width=\textwidth]{ecai-template/figures/graphnodesshap.png}
        \caption{GAT (Nodes)}
        \label{subfig:GATSHAP}
    \end{subfigure}
    \vspace{0.35cm}
    \caption{The top panels (a,b) show SHAP edge‐importance heatmaps for LR (left) and GAT (right). The bottom panels (c,d) depict nodes in Atlas coordinates space colored red (top 10), green (bottom 10), and yellow (others). The regions in (c) and (d) correspond to the region names in Table \ref{tab:regions_best_worst_overallheader} (NB// Node numbers here from 0-89). \\\vspace{1cm} }
    
    \label{fig:combined_importance}
\end{figure*}

% \begin{figure*}[htbp]
%     \centering
%     \begin{subfigure}[t]{0.48\textwidth}
%         \includegraphics[width=\textwidth]{ecai-template/figures/matrixnodesSHAP.png}
%         \label{subfig:LRSHAP}
%     \end{subfigure}%
%     \hfill
%     \begin{subfigure}[t]{0.48\textwidth}
%         \includegraphics[width=\textwidth]{ecai-template/figures/graphnodesshap.png}
%         \label{subfig:GATSHAP}
%     \end{subfigure}
%     \caption{Node numbers correspond to table
%     \ref{tab:regions_best_worst} \vspace{1cm}}
%     \label{fig:brainconnectivity_combined_large2}
% \end{figure*}




% \begin{figure}[htbp]
%     \centering
%     \begin{subfigure}[t]{0.48\columnwidth} 
%         \includegraphics[width=\columnwidth]{ecai-template/figures/matrixheatmap.png}
%         %\caption{Brain from above.}
%         \label{subfig:LRheatmap}
%     \end{subfigure}
%     \begin{subfigure}[t]{0.48\columnwidth}
%         \includegraphics[width=\columnwidth]{ecai-template/figures/graphheatmap.png}
%         %\caption{Brain from the front.}
%         \label{subfig:GATheatmap}
%     \end{subfigure}
%     \newline
%     \caption{ \vspace{1cm}}
%     \label{fig:brainconnectivity_combined}
% \end{figure}

% \begin{figure}[htbp]
%     \centering
%     \begin{subfigure}[t]{0.48\columnwidth} 
%         \includegraphics[width=\columnwidth]{ecai-template/figures/matrixnodesSHAP.png}
%         %\caption{Brain from above.}
%         \label{subfig:LRSHAP}
%     \end{subfigure}
%     \begin{subfigure}[t]{0.48\columnwidth}
%         \includegraphics[width=\columnwidth]{ecai-template/figures/graphnodesshap.png}
%         %\caption{Brain from the front.}
%         \label{subfig:GATSHAP}
%     \end{subfigure}
%     \newline
%     \caption{Node numbers correspond to table
%     \ref{tab:regions_best_worst} \vspace{1cm}}
%     \label{fig:brainconnectivity_combined}
% \end{figure}
\begin{figure}[h]
    \centering
        \includegraphics[width=\columnwidth]{ecai-template/figures/shap_summary_plot.png}

    \caption{SHAP summary plot visualizing each edge connection’s impact on the best model output (LR).}
     \vspace{1cm}
    \label{fig:connectionsshap}
\end{figure}


\section{Discussion}

From Table \ref{tab:classification_spatial} and Table \ref{tab:results_fuzzy} we see that LR is the best performing algorithm overall. Most models do well on average, but SVM, MLP and RF have trouble with the preterm class. Among the GNNs, GCN achieves moderate performance across both classes. GAT outperforms GCN, likely because its attention mechanism highlights the most informative graph connections. Adding spatial coordinates improves LR performance, as the atlas location data provides useful information, but it does not benefit the GNN models because it introduces additional complexity causing overfitting in some variants. For example, we note that GIN overfits, suggesting that more complex GNN models do not necessarily outperform simpler models like LR in this context. However, future work could broaden the scope through further experimentation with alternative architectures. One prior study applied SVM to functional brain data and reported an accuracy of 84\% \cite{smyser2016prediction}. Moreover, incorporating fuzzy logic into LR further enhances its performance by allowing smoother transitions around the decision threshold and better handling uncertainty in feature values. In our experiments, SVM achieved a comparable accuracy of approximately 86\% but was nonetheless outperformed by other methods. Although we focus on structural rather than functional development, we observed that SVM exhibited a pronounced bias toward the majority class, which is an important limitation given the relative scarcity of data in this domain.

%In computing these models, we are able to compare with other suggested models from literature \cite{cui2022braingb}\cite{taoudi2022predicting}, and show that our proposed method is better suited to the task of classifying preterm and term infants based on their structural brain connectivity. 

In order to further understand these results, we consider the main regions and connections that contribute to the classification. We note that while SHAP is not enough to confirm any causal relationships, there are a number of interesting findings which are supported by biological literature.
From SHAP analysis of LR and GAT models, the regions indicated are predominantly from the deep brain structures like putamen and thalamus, confirming previous findings \cite{gilmore2018imaging,rea2015essential}. Additionally, it demonstrates that the models consistently attend to specific regions, suggesting notable differences between preterm and term brain connectivity. Many of the connections which were highlighted as important, were between nodes which were also identified as important. Table \ref{tab:regions_best_worst_overallheader} demonstrates that there are regions consistently attended to regardless of LR or GAT model, but also indicates that the regions that are most consistent across models are the ones least influential which provides valuable insights into the effect of preterm birth on different brain regions.
Moreover, regions least important in distinguishing between the classes include regions which are generally understood to be well developed early before birth, for example Heschl’s gyrus, occipital lobe areas, and the temporal pole, confirming \citet{gilmore2018imaging}. These are associated with vision, hearing and other senses \cite{rea2015essential}. While regions most influential in the classification are understood to be present early in development, connections between these regions undergo significant development closer to the time of birth \cite{taymourtash2023fetal}. These regions have also been suggested as contributors in conditions such as Autism Spectrum Disorder (ASD)\cite{schuetze2016morphological}. The thalamus is one key region highlighted in our results. Neuroscientific studies have suggested that connective differences in the thalamus are linked to epilepsy \cite{marcuse2025thalamus}, and it is well established that the risk for epilepsy is increased by preterm birth \cite{ream2018neurologic,ball2012effect}. Another region which was particularly indicated in the LR model was cingulum, which has previously been implicated in developmental conditions following preterm birth \cite{lee2021cingulum}. Moreover, Figure \ref{fig:connectionsshap} illustrates that the connections between certain regions are influential in the classification, and demonstrates that higher feature values on certain edges can indicate either preterm or term depending on which regions they connect. For example this suggests that for higher feature values of the connection between Frontal\_Sup\_L and Frontal\_Sup\_Medial\_L, the model predicts cases of term infants. For the same edge with lower feature values, the model is more likely to classify as preterm. Many of the regions indicated in Figure \ref{tab:regions_best_worst_overallheader} are also indicated in Figure \ref{fig:combined_importance}, suggesting that not only are the nodes important, but suggesting that important nodes also connect to other important nodes \cite{bauer2017nonlinear,karolis2016reinforcement}.


% brain development is hindered by preterm birth, as teh scale of increase is altered with the changed enviroment, suggests an important order of development in the womb. thye also suggest that this could be helpful knowledge for predicting future problem for the children \cite{kapellou2006abnormal}. - it is important to consider individual differences - fuzzy

Including fuzzy logic in the model is important, as not only does it improve the accuracy and precision of the classification, but it allows for individual differences. Individuals close to the term cut off of 37 are difficult to classify, which may suggest that some are more similar to the term class than others of the same GA. This is important for further studies to take into account, and may also explain why it is difficult to predict pediatric outcomes for this particular group of infants \cite{engle2006recommendation,karnati2020late}.

Future work could consider alternative explainable models, and compare the insights with those we presented from SHAP, as well as comparisons to the neuroscientific knowledge. Moreover, our study is limited by data scarcity, so future studies should aim to incorporate any further data which becomes available. While we addressed the issue of class inbalance through undersampling, further work could explore alternative methods for counteracting this.   



\section{Conclusion}


In this paper, we compare matrix‑based and graph‑based classifiers for distinguishing preterm from term infant brain connectomes and show that adding spatial atlas coordinates improves model performance. We apply a fuzzy‑logic boundary around 37 weeks’ GA to the best LR and GAT models to smooth the decision threshold, raising accuracy from 0.8857 to 0.9333 with spatial coordinates and to 0.9619 when adding fuzzy logic, which lets the model learn gradual maturational changes. Among graph neural networks, GAT outperforms GCN but still falls short of LR, demonstrating that a simple linear model can separate both classes without added complexity.

Using SHAP edge and node importance scores, we identify a consistent set of deep‑brain regions (thalamus, putamen, cingulum) and their connections as the primary classification drivers, aligning with known neurodevelopmental findings on preterm risk. This interpretable mapping suggests potential biomarkers for early developmental screening. Future work can expand the cohort and treat gestational age as a continuous variable to refine sensitivity in the late‑preterm window and support more personalized assessments.


% KB - make this one section
% \subsection{Advantages}
% - comparison with the statistical methods/graph theoretical methods
% - eg Rich club results
% important nodes
% A key benefit of this analysis is the ability to integrate fuzzy logic. The assumption that full term birth is 37 weeks
%  This varys by country, and is likely to vary based on individual differences. If infants born close the 37 mark are hard to categorise it indicates that the 37  boundary is too fixed. 
% also calculating estimated conception date and due date can be an inexact process anyway. And the 37 may indicate that this is an average length of full term pregnancy rather than a hard boundary. 
% further research into this classification should consider usiing continuous age, or focus on a better way to classify?

% \subsection{Further suggestions}
% - continuous age 

%%%%%%%%%%%%%%%%%%%%%%%%%%%%%%%%%%%%%%%%%%%%%%%%%%%%%%%%%%%%%%%%%%%%%%%%
% \section{Citations and references}



%%%%%%%%%%%%%%%%%%%%%%%%%%%%%%%%%%%%%%%%%%%%%%%%%%%%%%%%%%%%%%%%%%%%%%%%





        
%%% Use this environment to include acknowledgements (optional).
%%% This will be omitted in doubleblind mode.

\begin{ack}
By using the \texttt{ack} environment to insert your (optional) 
acknowledgements, you can ensure that the text is suppressed whenever 
you use the \texttt{doubleblind} option. In the final version, 
acknowledgements may be included on the extra page intended for references.
\end{ack}

%%%%%%%%%%%%%%%%%%%%%%%%%%%%%%%%%%%%%%%%%%%%%%%%%%%%%%%%%%%%%%%%%%%%%%%%

%%% Use this command to include your bibliography file.

\clearpage
\bibliography{mybibfile}

\appendix
\section{Histogram of Gestational Age}

\begin{figure}[h]
    \centering
        \includegraphics[width=\columnwidth]{ecai-template/figures/GA.png}

    \caption{Gestational age histogram (including cases between 35–37 weeks). Demonstrates class imbalance.}
     % \vspace{1cm}
    \label{fig:connectionsshap}
\end{figure}
% \section{List of Nodes}


% \section{SHAP values for combined node metrics}

\end{document}
%%%%%%%%%%%%%%%%%%%%%%%%%%%%%%%%%%%%%%%%%%%%%%%%%%%%%%%%%%%%%%%%%%%%%%
